\chapter{Fundamentos teóricos}
\label{ch:fundamentos}
% Definiciones formales.

Este capítulo fija las definiciones formales que el Capítulo~\ref{ch:estado-del-arte} maneja en su versión conceptual. La Sección~\ref{sec:notacion} reúne la notación común a modo de glosario de símbolos. Sobre ella se establecen el descenso por gradiente estocástico y la geometría de los gradientes por ejemplo (Secciones~\ref{sec:sgd} y~\ref{sec:geometria}), y a continuación se definen las dos familias de métricas, la de alineación y la de variabilidad (Secciones~\ref{sec:def-alineacion} y~\ref{sec:def-variabilidad}). La Sección~\ref{sec:eficiencia} define los indicadores de eficiencia que actúan como variables a predecir, y la Sección~\ref{sec:fund-conclusiones} recapitula lo establecido. El capítulo define los objetos y, cuando la definición lo exige, anticipa qué agregado concreto se emplea. El protocolo de medición (el \emph{batch} de medición, las ventanas, el suavizado) y las hipótesis con signo asociadas a cada métrica se detallan en el Capítulo~\ref{ch:metodologia}.

\section{Notación}
\label{sec:notacion}

La Tabla~\ref{tab:notacion} recoge los símbolos que se emplean a lo largo de la memoria, con el lugar donde cada uno se introduce y define. Funciona como glosario de consulta. Cada símbolo se define una sola vez, en la sección donde aparece por primera vez, y esta tabla es el punto desde el que localizar esa definición sin releer el capítulo entero.

\begin{table}[!ht]
\centering
\small
\caption{Notación común de la memoria.}
\label{tab:notacion}
\begin{tabular}{cll}
\toprule
\textbf{Símbolo} & \textbf{Significado} & \textbf{Se introduce en} \\
\midrule
$(x_i, y_i)$          & ejemplo de entrenamiento $i$-ésimo                & Sec.~\ref{sec:sgd} \\
$n$                   & número de ejemplos de entrenamiento               & Sec.~\ref{sec:sgd} \\
$f_\theta$            & red neuronal con parámetros $\theta$              & Sec.~\ref{sec:sgd} \\
$P$                   & número de parámetros del modelo                   & Sec.~\ref{sec:sgd} \\
$\ell_i(\theta)$      & \emph{loss} del ejemplo $i$                              & Sec.~\ref{sec:sgd} \\
$\mathcal{L}(\theta)$ & riesgo empírico                                   & Sec.~\ref{sec:sgd} \\
$\eta_t$              & \emph{learning rate} en el paso $t$                & Sec.~\ref{sec:sgd} \\
$B$                   & tamaño del \emph{minibatch}                              & Sec.~\ref{sec:sgd} \\
$\hat{g}$             & gradiente del \emph{minibatch}                           & Sec.~\ref{sec:sgd} \\
$\Sigma$              & covarianza del gradiente por ejemplo              & Sec.~\ref{sec:sgd} \\
$g_i$                 & gradiente por ejemplo                             & Sec.~\ref{sec:geometria} \\
$M$                   & tamaño del \emph{batch} de medición                   & Sec.~\ref{sec:geometria} \\
$\mathbf{G}$          & matriz de Gram de los $g_i$                       & Sec.~\ref{sec:geometria} \\
$s$                   & submuestras disjuntas en la disparidad, $s = 5$ de 51 ejemplos & Sec.~\ref{sec:def-alineacion} \\
$\bar{g}^{(a)}$       & gradiente medio de la submuestra $a$              & Sec.~\ref{sec:def-alineacion} \\
$W$                   & pesos de la última capa lineal                    & Sec.~\ref{sec:def-alineacion} \\
$K$                   & submuestras disjuntas en la NGV, $K = 10$ de 25 ejemplos & Sec.~\ref{sec:def-variabilidad} \\
$\bar{g}$             & gradiente medio del \emph{batch} de medición          & Sec.~\ref{sec:def-variabilidad} \\
$S$, $Q$              & suma de los $g_i$ y suma de sus normas al cuadrado & Sec.~\ref{sec:def-variabilidad} \\
$T$                   & presupuesto de \emph{epochs}                             & Sec.~\ref{sec:eficiencia} \\
$\tau$                & umbral de \emph{accuracy} de validación de la celda      & Sec.~\ref{sec:eficiencia} \\
$t^\star$             & \emph{epochs} hasta el umbral                            & Sec.~\ref{sec:eficiencia} \\
$f$                   & fracción temprana del presupuesto                 & Cap.~\ref{ch:metodologia} \\
\bottomrule
\end{tabular}
\end{table}

% La tabla de notación no cabe en lo que la apertura del capítulo deja libre en la página, ni
% siquiera a cuerpo menor, así que sin este corte LaTeX la empuja dentro de \ref{sec:sgd},
% fuera de la sección que la presenta. El precio son dos páginas más en el total.
\clearpage

\section{Descenso por gradiente estocástico}
\label{sec:sgd}

Sea un conjunto de entrenamiento de $n$ ejemplos $\{(x_i, y_i)\}_{i=1}^n$, una red $f_\theta$ con parámetros $\theta \in \mathbb{R}^P$ y una función de \emph{loss} por ejemplo $\ell_i(\theta) = \ell(f_\theta(x_i), y_i)$. El objetivo del entrenamiento es minimizar el riesgo empírico $\mathcal{L}(\theta) = \tfrac{1}{n}\sum_{i=1}^n \ell_i(\theta)$. El descenso por gradiente estocástico con \emph{minibatch} no evalúa el gradiente exacto $\nabla \mathcal{L}(\theta)$, sino una estimación sobre un subconjunto $I_t$ de $B$ ejemplos extraídos al azar en cada paso:

\begin{equation}
\theta_{t+1} = \theta_t - \eta_t\, \hat{g}_t,
\qquad
\hat{g}_t = \frac{1}{B}\sum_{i \in I_t} \nabla_\theta \ell_i(\theta_t),
\label{eq:sgd}
\end{equation}

donde $\eta_t$ es el \emph{learning rate}. Si el \emph{minibatch} se muestrea de manera uniforme, el estimador es \emph{insesgado}, $\mathbb{E}[\hat{g}_t] = \nabla \mathcal{L}(\theta_t)$, lo que garantiza que SGD optimiza el \emph{loss} correcto y no una distorsión sistemática de él. Su dispersión, en cambio, no es nula. Con ejemplos independientes, la traza de su covarianza vale

\begin{equation}
\operatorname{tr}\operatorname{Cov}[\hat{g}_t] = \frac{\operatorname{tr}(\Sigma)}{B},
\qquad
\Sigma = \operatorname{Cov}_i\!\big[\nabla_\theta \ell_i(\theta_t)\big],
\label{eq:sgd-var}
\end{equation}

con $\Sigma$ la covarianza del gradiente por ejemplo. Esta dispersión escala como $1/B$ pero persiste incluso en el óptimo, donde $\nabla \mathcal{L} = 0$ y aun así $\operatorname{tr}\operatorname{Cov}[\hat{g}] > 0$, porque cada \emph{minibatch} apunta a un objetivo ligeramente distinto. Esa dispersión residual es el suelo de ruido que obliga a reducir progresivamente el \emph{learning rate} para converger, en lugar de orbitar alrededor del mínimo, y es la cantidad que cuantifica la familia de variabilidad.

Los dos optimizadores del estudio parten de ese mismo gradiente y lo transforman de forma distinta. SGD con \emph{momentum} acumula una media móvil de los gradientes y da el paso con ella, y Adam~\cite{kingma2015adam} mantiene dos medias móviles, una del gradiente y otra de su cuadrado por coordenada, y divide la primera por la raíz de la segunda:

\begin{equation}
\text{SGD:}\quad m_t = \mu\, m_{t-1} + \hat{g}_t,\ \ \theta_{t+1} = \theta_t - \eta\, m_t;
\qquad
\text{Adam:}\quad \theta_{t+1} = \theta_t - \eta\, \frac{\hat{m}_t}{\sqrt{\hat{v}_t} + \varepsilon},
\label{eq:momentum-adam}
\end{equation}

donde $\hat{m}_t$ y $\hat{v}_t$ son las dos medias móviles de Adam corregidas por su sesgo inicial. Con $\mu = 0{,}9$ y un gradiente estable, $m_t$ tiende a $\hat{g}/(1 - \mu) = 10\,\hat{g}$, así que el paso efectivo de SGD es unas diez veces el de la Ecuación~\ref{eq:sgd} y sigue siendo proporcional al gradiente. El paso de Adam, al dividir cada coordenada por su propia escala, mide del orden de $\eta$ por coordenada sea cual sea la magnitud del gradiente. Esa diferencia es la razón práctica por la que la rejilla de \emph{learning rates} de Adam se sitúa una década por debajo de la de SGD en la Sección~\ref{sec:matriz}. El argumento no es una igualdad de magnitudes, porque los dos pasos dependen del gradiente de manera distinta; se apoya en los valores por defecto canónicos, $10^{-2}$ para SGD con \emph{momentum} y $10^{-3}$ para Adam, y en una comprobación sobre la matriz: sobre las ocho posiciones de cada rejilla se quedan en el azar 73 entrenamientos con SGD y 81 con Adam, de 480 cada uno, de modo que las dos listas cubren un tramo comparable de su rango útil.

Dos resúmenes estadísticos del gradiente por ejemplo $g$ (que la Sección~\ref{sec:geometria} define formalmente) organizan el resto del capítulo. El \emph{primer momento} $\mathbb{E}[g]$ describe la dirección consistente del gradiente y es lo que estiman el término de inercia (\emph{momentum}) de SGD y la media móvil $m_t$ de Adam~\cite{kingma2015adam}. El \emph{segundo momento}, a través de $\mathbb{E}[\|g\|^2]$ o de la traza de $\Sigma$, describe su magnitud y su dispersión, y es lo que recogen la media móvil $v_t$ de Adam y la escala de ruido del gradiente. Las dos familias de métricas leen, respectivamente, la dirección relativa de los gradientes por ejemplo (alineación) y su dispersión (variabilidad).

Esa lectura fija además una decisión de diseño central del trabajo: todas las métricas se calculan sobre el gradiente bruto $\nabla \ell$ y nunca sobre el paso ya preacondicionado por el optimizador. El motivo es que Adam divide el gradiente por una estimación de su propio segundo momento, que es justamente lo que mide la familia de variabilidad. Medir sobre el paso haría que una misma red, con el mismo gradiente, arrojase valores distintos según el optimizador, y la métrica pasaría a describir en parte la mecánica de este en lugar del aprendizaje. Sobre el gradiente bruto, en cambio, las ocho métricas son la misma cantidad bajo SGD y bajo Adam, que es lo que hace legítima la comparación entre ambos en el Capítulo~\ref{ch:metodologia}.

\section{Geometría del gradiente}
\label{sec:geometria}

El objeto sobre el que se construyen casi todas las métricas es el gradiente por ejemplo

\begin{equation}
g_i = \nabla_\theta\, \ell(f_\theta(x_i), y_i) \in \mathbb{R}^P,
\label{eq:per-sample}
\end{equation}

el gradiente del \emph{loss} de un único ejemplo, antes de promediar. El gradiente del \emph{minibatch} es su media $\hat{g} = \tfrac{1}{B}\sum_i g_i$, pero esa media descarta la estructura estadística que las métricas explotan. La distribución de los $g_i$ (su covarianza, su curtosis) describe el ruido, y sus productos internos por pares describen si los ejemplos cooperan o se cancelan. Para medirlos sin contaminar la dinámica de entrenamiento, todas las métricas se evalúan siempre sobre el mismo \emph{batch} de medición. Se denomina así al subconjunto de $M$ ejemplos, extraído del conjunto de entrenamiento, que se selecciona una única vez al comenzar el entrenamiento y se reutiliza sin cambios en todas las mediciones posteriores. No es una partición de datos adicional junto a las de entrenamiento, validación y test, sino un instrumento de medida. Los ejemplos que lo forman siguen participando en el entrenamiento con normalidad. Se extrae por permutación aleatoria del conjunto de entrenamiento, sin estratificar por clase, así que el número de ejemplos de cada clase que contiene varía con la \emph{seed}. Con $M = 256$ ejemplos y $C$ clases, cabe esperar $\binom{M}{2}/C$ pares de la misma clase: unos 3.264 en MNIST y CIFAR-10, unos 326 en CIFAR-100 y unos 163 en Tiny-ImageNet, de los 32.640 pares posibles. La Sección~\ref{sec:def-alineacion} vuelve sobre ello, porque la \emph{stiffness} intra-clase descansa en esos pares.

Que permanezca fijo es precisamente lo que permite comparar una métrica consigo misma a lo largo del entrenamiento. Si en cada medición se emplearan ejemplos distintos, una variación de la métrica entre dos instantes admitiría dos explicaciones incompatibles, un cambio en el modelo o un remuestreo afortunado, y no habría forma de separarlas. Al congelar los ejemplos, la única fuente de variación que queda es el modelo.

Todas las métricas se calculan con la red en modo de evaluación, esto es, con las capas de \emph{batch normalization} usando las estadísticas acumuladas durante el entrenamiento en lugar de las del lote. En la red \emph{fully connected} y en la convolucional no hay ninguna capa de ese tipo y el modo no cambia nada. En la ResNet-18, la única de las tres que la lleva, es la única forma de definir un gradiente por ejemplo, porque en modo de entrenamiento la salida de cada ejemplo depende del resto del lote. A cambio, el gradiente que se mide es el de la red con sus estadísticas acumuladas, que en las primeras \emph{epochs} no coinciden aún con las del lote, y no exactamente el que desciende el optimizador. Es una diferencia de objeto medido que hay que tener delante al comparar la ResNet-18 con las otras dos arquitecturas.

Las dos operaciones elementales son la matriz de Gram y la similitud coseno. La matriz de Gram recoge todos los productos internos por pares, $\mathbf{G}_{ij} = \langle g_i, g_j\rangle$, y la similitud coseno los normaliza por magnitud:

\begin{equation}
\cos(g_i, g_j) = \frac{\langle g_i, g_j\rangle}{\|g_i\|\,\|g_j\|} \in [-1, 1].
\label{eq:cosine}
\end{equation}

El coseno vale $1$ si los dos gradientes apuntan en la misma dirección, $0$ si son ortogonales y $-1$ si son opuestos. La normalización por norma importa porque lo relevante para la cooperación entre ejemplos es la dirección relativa, no la escala individual de cada gradiente.

Conviene detenerse en un hecho de alta dimensión, porque fija la escala con la que hay que leer cualquier coseno entre gradientes. Se llama \emph{isótropo} a un vector aleatorio cuya distribución no privilegia ninguna dirección del espacio, esto es, que es invariante frente a rotaciones. El caso canónico es un vector gaussiano con covarianza proporcional a la identidad o, lo que es lo mismo una vez normalizado, un punto uniforme sobre la esfera unidad. Si $u$ y $v$ son dos vectores isótropos independientes en $\mathbb{R}^P$, su coseno tiene media nula y varianza exactamente $1/P$, de modo que su magnitud típica es del orden de $1/\sqrt{P}$. No se trata solo de que la media sea cero. La distribución se concentra exponencialmente en torno a ese valor, con $\Pr[\,|\cos(u,v)| \ge t\,]$ decayendo como $\exp(-P t^2 / 2)$, que es el fenómeno de concentración de la medida sobre la esfera~\cite{vershynin2018high}. Dicho de otro modo, en dimensión alta dos direcciones tomadas al azar son casi ortogonales, y lo son tanto más cuanto mayor es $P$~\cite{blum2020foundations}.

Ese cálculo fija la escala del ruido de muestreo, no la referencia contra la que leer una métrica. En los modelos de este trabajo $P$ va desde unos $1{,}9 \times 10^{4}$ parámetros en la CNN pequeña hasta más de $10^{7}$ en ResNet-18 y en la red \emph{fully connected} sobre Tiny-ImageNet, lo que sitúa el coseno entre dos direcciones sin relación alguna entre $7 \times 10^{-3}$ y $3 \times 10^{-4}$. Los gradientes por ejemplo de una red no son direcciones sin relación, ni siquiera en la inicialización. Con entropía cruzada, el gradiente respecto a la última capa es $(p - y)\,z^\top$, donde $p$ son las probabilidades predichas, $y$ la etiqueta en codificación \emph{one-hot} y $z$ la entrada de esa capa. Dos ejemplos de la misma clase comparten el signo del factor $(p - y)$ y, con activaciones ReLU no negativas, sus gradientes forman un ángulo agudo antes de haber aprendido nada, mientras que dos ejemplos de clases distintas tienden a oponerse. La \emph{stiffness} intra-clase es positiva y la inter-clase negativa por estructura, y la referencia honesta de cada métrica de alineación es su valor en la inicialización. La matriz no lo registra, porque la primera medición se hace al final de la primera \emph{epoch}. Los valores registrados muestran esa estructura: sobre las 33.600 filas de la matriz, el coseno medio entre todos los pares del \emph{batch} de medición tiene mediana $0{,}0019$ en la CNN, $0{,}0005$ en la red \emph{fully connected} y $0{,}009$ en la ResNet-18, es decir, en el suelo del ruido en las dos primeras y unas treinta veces por encima en la tercera, mientras que el coseno entre pares de la misma clase tiene mediana entre $0{,}30$ y $0{,}51$ en las tres y el de pares de clases distintas se queda en cero. Lo que garantiza el argumento de concentración es que una desviación como esa es un hecho de la red y no una fluctuación del muestreo. El argumento alcanza también a los estadísticos de extremo. Al recorrer los $\binom{M}{2}$ pares del \emph{batch} de medición, el coseno más negativo esperable entre direcciones sin relación solo alcanza en magnitud el orden de $\sqrt{\log M / P}$, de manera que tampoco un valor apreciable de \emph{gradient confusion} puede atribuirse al azar del muestreo.

\section{Métricas de alineación}
\label{sec:def-alineacion}

La familia de alineación cuantifica en qué medida los gradientes de ejemplos distintos apuntan hacia el mismo sitio. Su marco conceptual es la descomposición de la norma del gradiente del \emph{minibatch}~\cite{chatterjee2020coherent},

\begin{equation}
\Big\|\sum_i g_i\Big\|^2 = \sum_i \|g_i\|^2 + \sum_{i \neq j} \langle g_i, g_j\rangle,
\label{eq:cgh}
\end{equation}

donde el primer término suma el tamaño al cuadrado del gradiente de cada ejemplo, esto es, lo que cada uno aporta por su cuenta, y el segundo suma el producto interno de cada par de gradientes distintos, positivo cuando los dos apuntan al mismo lado y negativo cuando se oponen. En ese segundo término vive la coherencia, y las métricas siguientes son distintos agregadores de él.

\paragraph{m-coherencia.} Reduce la coherencia a un único escalar sobre el \emph{batch} de medición de $M$
ejemplos~\cite{chatterjee2020making}:

\begin{equation}
\alpha = \frac{\big\|\sum_{i=1}^{M} g_i\big\|^2}{\sum_{i=1}^{M} \|g_i\|^2} \in [0, M].
\label{eq:mcoherence}
\end{equation}

Vale $M$ cuando los gradientes son idénticos, en dirección y en norma, y $1$ en el \emph{límite ortogonal}, donde son mutuamente perpendiculares, cualquiera que sea su norma. Por debajo de $1$ los gradientes están anticorrelados. Su lectura intuitiva es el número medio de ejemplos que se benefician de un paso de descenso obtenido a partir del gradiente de un ejemplo cualquiera. Una alineación mayor se asocia con un aprendizaje más coherente, de modo que se espera que una $\alpha$ alta acompañe a un entrenamiento eficiente.

\paragraph{Stiffness.} Mide el coseno entre gradientes de pares de ejemplos, desagregado
por clases~\cite{fort2019stiffness}. Sobre el \emph{batch} de medición se definen una variante de coseno y una de signo,

\begin{equation}
S_{\cos} = \mathbb{E}_{i \neq j}\!\big[\cos(g_i, g_j)\big],
\qquad
S_{\operatorname{sign}} = \mathbb{E}_{i \neq j}\!\big[\operatorname{sign}\langle g_i, g_j\rangle\big],
\label{eq:stiffness}
\end{equation}

cada una restringida a pares de la misma clase ($y_i = y_j$, \emph{stiffness} \emph{intra-clase}) o de clases distintas ($y_i \neq y_j$, \emph{inter-clase}). Una \emph{stiffness} intra-clase alta indica que la red ha aprendido características compartidas y que cada paso de gradiente sobre un ejemplo beneficia al resto de su clase. La columna que representa a la métrica es el coseno intra-clase, según fija la Sección~\ref{sec:variables}, y la media se toma sobre todos los pares mezclados, no como media de una matriz por clases como en el artículo. Con el \emph{batch} de medición sin estratificar de la Sección~\ref{sec:geometria}, ese coseno descansa en unos 163 pares en Tiny-ImageNet y unos 326 en CIFAR-100, y su composición cambia con la \emph{seed}. Fort et al. afirman que la señal predictiva es más fuerte en la fase temprana, porque la \emph{stiffness} decae hacia cero cuando comienza el sobreajuste; es una afirmación del artículo y no un hecho medido aquí.

\paragraph{Gradient confusion.} Adopta el punto de vista del peor caso y toma el coseno más
negativo entre cualquier par de gradientes del \emph{batch} de medición~\cite{sankararaman2020impact}. En su forma normalizada,

\begin{equation}
\hat{\eta} = -\min_{i \neq j} \cos(g_i, g_j).
\label{eq:confusion}
\end{equation}

Un $\hat{\eta}$ alto señala que existe al menos un par de ejemplos cuyos gradientes se oponen con fuerza, lo que frena el avance neto de SGD. Un $\hat{\eta}$ cercano a cero indica gradientes que no se contradicen. Por ser un estadístico de extremo es ruidoso, así que junto al mínimo se registran estadísticos de la distribución de cosenos: su mediana, su percentil 5 y la fracción de pares negativos. El artículo de origen la estima sobre 100 pares de gradientes de \emph{minibatch} de 128 ejemplos al final de cada \emph{epoch}. Aquí es el mínimo sobre los 32.640 pares de gradientes por ejemplo del \emph{batch} de medición, un extremo sobre muchos más pares y de otra distribución, así que su magnitud no es comparable con las cifras del artículo, aunque su orden dentro de una celda sí lo es. La columna titular se mantiene en el mínimo por fidelidad a la definición original. Se predice que una confusión menor acompaña a un entrenamiento más rápido.

\paragraph{Gradient disparity.} Sustituye el coseno por la distancia entre los gradientes
medios de dos submuestras independientes, derivada de una cota de tipo PAC-Bayes~\cite{forouzesh2021disparity}. Si el \emph{batch} de medición se parte en $s$ submuestras disjuntas y $\bar{g}^{(a)}$ denota el gradiente medio de la submuestra $a$,

\begin{equation}
\overline{\mathcal{D}} = \binom{s}{2}^{-1} \sum_{a < b} \big\|\bar{g}^{(a)} - \bar{g}^{(b)}\big\|_2.
\label{eq:disparity}
\end{equation}

Es una magnitud $\ell_2$ \emph{sin normalizar}, no una similitud coseno, porque la cota teórica del trabajo depende de $\|\bar{g}^{(a)} - \bar{g}^{(b)}\|_2^2$ y dividir por la norma rompería la conexión con la teoría. Aquí $s = 5$ y cada submuestra tiene 51 ejemplos. Esa forma tiene una consecuencia que la definición no deja ver. Si las submuestras son disjuntas y de tamaño $b$, la esperanza de $\|\bar{g}^{(a)} - \bar{g}^{(b)}\|^2$ es $2\operatorname{tr}(\Sigma)/b$, y en dimensión alta la norma se concentra en su esperanza, así que la disparidad es, salvo el ruido de muestreo, la raíz de $2\operatorname{tr}(\Sigma)/51$: la dispersión del gradiente por ejemplo sin normalizar por la señal. Sobre las 33.600 filas de la matriz, el cociente entre el cuadrado de la disparidad y la traza de la covarianza tiene mediana $0{,}038$, con la mitad central entre $0{,}036$ y $0{,}040$, frente al $2/51 = 0{,}039$ del álgebra, y su correlación de Spearman con esa traza dentro de un mismo entrenamiento tiene mediana $0{,}99$. Por eso, desde el 3 de septiembre de 2026, esta métrica cuenta en la familia de variabilidad y no en la de alineación, como recoge la Sección~\ref{sec:protocolo-objetivos}. Es además la única de las ocho que no es invariante a la escala del gradiente, de modo que dentro de una celda recoge en parte cuánto mide el gradiente en ese momento. Una disparidad baja indica submuestras coherentes, y se predice como señal de eficiencia, con la salvedad de que su magnitud depende del tamaño de la submuestra, 51 aquí, y solo es comparable entre entrenamientos que lo fijen.

\paragraph{Alineación gradiente-peso (GWA).} Mide el coseno entre el gradiente por ejemplo
y los pesos de la última capa lineal de la red, la que convierte la representación aprendida en una puntuación por clase, con el gradiente restringido también a esa capa~\cite{hoelzl2025gradient}. Con $W$ los pesos de esa capa, sin el sesgo, y $g_i^{(W)}$ el gradiente de $\ell_i$ respecto a ellos, para cada ejemplo del \emph{batch} de medición se toma $\gamma_i = \cos\!\big(g_i^{(W)}, W\big)$, y el agregado por \emph{epoch} corrige la media por el exceso de curtosis de la distribución de los $\gamma_i$:

\begin{equation}
\mathrm{GWA} = \frac{\bar{\gamma}}{\dfrac{\mu_4}{\mu_2^{2}} - 3 + \beta},
\qquad \beta = 1{,}2,
\label{eq:gwa}
\end{equation}

donde $\bar{\gamma}$ es la media de las $\gamma_i$, $\mu_2$ y $\mu_4$ sus momentos centrales de segundo y cuarto orden, y $\mu_4/\mu_2^2 - 3$ el exceso de curtosis (nulo para una gaussiana). El denominador descuenta las distribuciones de cola pesada, donde pocas muestras dominan la media.

Cada $\gamma_i$ tiene una lectura directa en términos de la red. Con entropía cruzada y logits $Wz + b$, el producto interno entre el gradiente de $W$ y la propia $W$ vale $\sum_c (p_c - y_c)\,(Wz)_c$, la diferencia entre el logit medio bajo las probabilidades predichas y el logit de la clase correcta, y $\gamma_i$ es esa cantidad dividida por las dos normas. Cuanto más domina el logit de la clase correcta, más negativo es $\gamma_i$. GWA es, por tanto, un margen normalizado medio de la red sobre los 256 ejemplos del \emph{batch} de medición, corregido por la forma de su distribución. Eso explica que en su artículo siga tan de cerca al \emph{accuracy} de test, porque es una cantidad que el propio paso hacia delante casi regala, y anticipa que dentro de una celda irá muy cerca del \emph{accuracy} de validación leído en la misma ventana. Los valores registrados lo confirman en el rango: la media $\bar{\gamma}$ vive entre $-0{,}18$ y $0{,}02$ en toda la matriz. El cociente de la Ecuación~\ref{eq:gwa}, en cambio, recorre de $-63$ a $131$, porque el denominador se hace negativo cuando el exceso de curtosis baja de $-1{,}2$, cosa que ocurre en el 1,8\,\% de las filas registradas y en ninguna de la ventana del 5\,\% que lee el análisis primario. En esas filas GWA cambia de signo sin que la media lo haga, y las ventanas exploratorias se leen con ese aviso. Hay además una salvedad de signo que condiciona toda lectura posterior de esta métrica. El trabajo de origen define el gradiente como $-\nabla\ell$ y concluye que una GWA \emph{alta} se asocia con mejor generalización. Como aquí se mide sobre el gradiente bruto $\nabla\ell$, por la decisión de la Sección~\ref{sec:sgd}, el signo queda invertido respecto al original. La predicción heredada, expresada en la convención de esta memoria, es que una GWA \emph{baja} acompañe a una mejor generalización. El Capítulo~\ref{ch:metodologia} contrasta el signo observado contra esta forma ya convertida, y no contra la del artículo.

Cada una de estas definiciones lleva implícita una predicción direccional: más coherencia, más \emph{stiffness} intra-clase, menos confusión y menos disparidad deberían acompañar a un entrenamiento más eficiente, y, con la salvedad de signo que se acaba de describir, una GWA más baja a una mejor generalización. Esas predicciones con signo se formalizan como hipótesis falsables en el Capítulo~\ref{ch:metodologia}.

\section{Métricas de variabilidad}
\label{sec:def-variabilidad}

La familia de variabilidad no mira la dirección relativa de los gradientes, sino su dispersión respecto a la señal. Sus tres métricas, y desde el 3 de septiembre de 2026 también la \emph{gradient disparity}, son facetas de una misma cantidad, la dispersión del gradiente por ejemplo medida contra su media.

\paragraph{Varianza normalizada (NGV).} Hace comparable la varianza del gradiente entre
problemas de distinta escala dividiéndola por la magnitud de la señal~\cite{faghri2020study}. Se estima sobre $K = 10$ submuestras disjuntas de 25 ejemplos del \emph{batch} de medición, cuyos gradientes medios aportan las $K$ observaciones sobre las que se toman la esperanza y la covarianza:

\begin{equation}
\mathrm{NGV} = \frac{\operatorname{tr}\!\big(\operatorname{Cov}(g)\big)}{\|\mathbb{E}[g]\|^2}.
\label{eq:ngv}
\end{equation}

Es el inverso de una razón señal-ruido. Por debajo de $1$ la señal del gradiente medio domina sobre el ruido, y por encima domina el ruido. Esta forma, de tipo traza sobre norma, es una adaptación deliberada de la versión literal del trabajo de origen ($\mathbb{V}[g_j]/\mathbb{E}[g_j^2]$ por coordenada), monótonamente equivalente. Una NGV menor indica menos ruido relativo, si bien empíricamente la métrica crece durante el entrenamiento en lugar de decrecer. La razón es elemental. Sin regularización, a medida que la red ajusta los ejemplos del \emph{batch} de medición el gradiente medio tiende a cero mientras la dispersión entre ejemplos no, así que el cociente sube. Dentro de una celda, el \emph{learning rate} decide cuánto ha avanzado la red cuando se mide, luego estas columnas son en buena parte cuánto ha bajado ya el \emph{loss}, leído con otra regla; es el mecanismo de la causa común que el Capítulo~\ref{ch:resultados} mide. El estimador tiene además un techo. El denominador se estima con la norma del gradiente medio de las $K$ submuestras, que está sesgada hacia arriba por $\operatorname{tr}(\Sigma)/K$, y por eso la NGV registrada no puede pasar de $K = 10$ aunque la señal sea nula; su mediana sobre la matriz es $8{,}1$, pegada a ese techo. Desde el 3 de septiembre de 2026 la NGV no entra en el contraste, porque es la escala de ruido de la Ecuación~\ref{eq:gns} estimada con diez submuestras en lugar de con los 256 ejemplos, y la escala de ruido la representa (Sección~\ref{sec:protocolo-objetivos}).

\paragraph{Escala de ruido del gradiente (GNS).} Predice el tamaño de \emph{batch} crítico, aquel
a partir del cual ampliar el \emph{batch} deja de reducir el número de pasos, sin necesidad de barrer ese hiperparámetro~\cite{mccandlish2018empirical}. Su forma simplificada es

\begin{equation}
\mathcal{B}_{\text{simple}} = \frac{\operatorname{tr}(\Sigma)}{\|\bar{g}\|^2},
\label{eq:gns}
\end{equation}

con $\Sigma$ la covarianza por ejemplo y $\bar{g}$ el gradiente medio. La versión exacta $\mathcal{B}_{\text{noise}} = \operatorname{tr}(H\Sigma)/(\bar{g}^\top H \bar{g})$ requiere la Hessiana $H$ y se descarta por coste. La simplificada asume $H \propto I$. Por construcción, $\mathcal{B}_{\text{simple}} \approx 25 \cdot \mathrm{NGV}$, donde 25 es el tamaño de las submuestras con las que se estima la NGV y no el tamaño de \emph{batch} del entrenamiento; la relación es exacta para las cantidades poblacionales y solo aproximada para los estimadores.

Estimada sobre el \emph{batch} de medición con gradientes por ejemplo, la escala de ruido guarda además una relación que no es aproximada sino una identidad algebraica, y que cruza la frontera entre las dos familias. Con $S = \sum_i g_i$ y $Q = \sum_i \|g_i\|^2$ se tiene $\|\bar{g}\|^2 = \|S\|^2/M^2$ y $\operatorname{tr}(\Sigma) = Q/M - \|S\|^2/M^2$, de donde

\begin{equation}
\mathcal{B}_{\text{simple}} = \frac{M}{\alpha} - 1,
\label{eq:gns-mcoh}
\end{equation}

con $\alpha$ la m-coherencia de la Ecuación~\ref{eq:mcoherence}. Ambas dependen solo del cociente $\|S\|^2/Q$, de modo que son la misma cantidad reparametrizada y su correlación de Spearman dentro de un mismo entrenamiento está forzada a $-1$ exacto. El piloto de calibración lo verificó a un error relativo de $8{,}8 \times 10^{-6}$, atribuible al redondeo en coma flotante. La consecuencia es que la escala de ruido y la m-coherencia no aportan evidencia independiente pese a pertenecer a familias distintas, y la NGV tampoco, porque es esta misma razón estimada con diez submuestras. Las tres son estimadores de $\alpha$, la razón entre $\|S\|^2$ y $Q$; solo GSNR agrega de otra manera, por coordenada. Desde el 3 de septiembre de 2026 la m-coherencia y la NGV salen del contraste y la escala de ruido las representa, como recoge la Sección~\ref{sec:protocolo-objetivos}, y la predicción de signo del artículo de la m-coherencia se lee sobre la escala de ruido con el signo cambiado, porque la Ecuación~\ref{eq:gns-mcoh} es decreciente en $\alpha$.

El estimador de la escala de ruido tiene un techo, y está cerca de él. Si el gradiente medio verdadero es cero, $\mathbb{E}\|S\|^2 = Q$, $\alpha$ tiende a $1$ y la escala de ruido estimada a $M - 1 = 255$. McCandlish et al. reportan escalas de $10^{3}$ a $10^{4}$ en CIFAR-10 e ImageNet, muy por encima de lo que 256 ejemplos pueden resolver, así que el estimador vive pegado a su techo: sobre las 33.600 filas de la matriz la mediana es $207$, el percentil 5 es $58$ y el percentil 95 es $287$, con la mitad de las filas por encima de $204$. Las magnitudes registradas no se pueden leer, por tanto, como el tamaño de \emph{batch} crítico que la métrica promete, y la transformación $M/\alpha - 1$ amplifica cerca de $\alpha = 1$ cualquier fluctuación de $\alpha$. Lo que sí conserva sentido es el orden entre entrenamientos de una misma celda, que es lo que el análisis usa, con menos potencia de la que tendría un estimador lejos de su techo. McCandlish et al. dan un estimador sin sesgo de la señal, $\|\bar{g}\|^2 - \operatorname{tr}(\hat{\Sigma})/M$, que puede salir negativo y declarar entonces que la señal no es detectable con 256 ejemplos; se puede calcular desde las columnas registradas y no se calculó.

\paragraph{Razón señal-ruido del gradiente por parámetro (GSNR).} Lleva la razón señal-ruido al nivel de
cada parámetro $\theta_j$~\cite{liu2020understanding}:

\begin{equation}
r(\theta_j) = \frac{\tilde{g}(\theta_j)^2}{\sigma^2(\theta_j)},
\qquad
\tilde{g}(\theta_j) = \mathbb{E}_i[g_{i,j}],
\quad
\sigma^2(\theta_j) = \operatorname{Var}_i[g_{i,j}],
\label{eq:gsnr}
\end{equation}

agregado por la media sobre los parámetros, nunca por la suma, que sería incomparable entre arquitecturas de distinto tamaño. El estimador usa la varianza insesgada sobre los $M$ ejemplos, descarta los parámetros cuyo gradiente es nulo en todo el \emph{batch} de medición, y bajo señal nula tiene un suelo de $1/M \approx 0{,}004$ por parámetro, porque el cuadrado de la media de $M$ valores de ruido vale en esperanza su varianza dividida por $M$. La media sobre millones de parámetros la domina la cola pesada de $r(\theta_j)$, y por eso se registran también la mediana y el percentil 95; la columna titular es la media por fidelidad a Liu et al. Es esencialmente el recíproco por parámetro de la NGV. Un GSNR alto significa que la media del gradiente por ejemplo domina a su varianza, y se asocia con un \emph{gap} de generalización pequeño. Las tres métricas miden, pues, la misma razón señal-ruido desde ángulos distintos. La redundancia entre ellas es algebraica y se resuelve por álgebra antes del contraste, como recoge la Sección~\ref{sec:protocolo-objetivos}: queda la escala de ruido como representante de $\alpha$, y GSNR aparte porque agrega por coordenada.

\section{Indicadores de eficiencia del entrenamiento}
\label{sec:eficiencia}

Las métricas anteriores son las variables predictoras. Las que siguen son las variables a predecir, esto es, lo que en este trabajo se entiende por eficiencia del entrenamiento. Se leen sobre la curva de validación a lo largo del presupuesto de \emph{epochs} $T$. El protocolo concreto, incluido el suavizado de la curva, la separación entrenamiento/validación/test y el tratamiento del censurado en el análisis, se especifica en el Capítulo~\ref{ch:metodologia}.

\paragraph{Epochs hasta el umbral (indicador primario).} El número de \emph{epochs} que tarda la
\emph{accuracy} de validación en alcanzar un nivel objetivo $\tau$ propio de cada conjunto de datos y arquitectura,

\begin{equation}
t^\star = \min\{\, t \le T : \mathrm{acc}_{\text{val}}(t) \ge \tau \,\}.
\label{eq:epochs-to-threshold}
\end{equation}

Si el \emph{accuracy} no cruza el umbral dentro del presupuesto, el valor queda \emph{censurado} y se anota como ausente, nunca como el valor del presupuesto. Los extremos de la rejilla de \emph{learning rates} no alcanzan el umbral por diseño. Descartar un entrenamiento lento pero vivo sesgaría la muestra, y por eso se conserva como censurado; distinto es un entrenamiento que nunca aprendió, que no es lento sino que está muerto, y la Sección~\ref{sec:protocolo-poblacion} lo deja fuera de la población con ese criterio. Un $t^\star$ menor indica un entrenamiento más rápido.

\paragraph{AUC del \emph{loss} de validación.} El área bajo la curva de \emph{loss} de validación
dentro del presupuesto, aproximada por la regla trapezoidal sobre el eje de \emph{epochs}. Resume toda la trayectoria en un escalar. Cuanto menor es el área, más rápido y más bajo desciende el \emph{loss}, siempre que la curva descienda. En el régimen de este estudio, sin \emph{weight decay} ni \emph{data augmentation}, el \emph{loss} de validación toca fondo pronto y vuelve a subir, y el área recoge entonces sobre todo esa subida: en las celdas de la red \emph{fully connected} y de la ResNet-18 fuera de MNIST el mínimo llega entre la \emph{epoch} 4 y la 8 de mediana, y entre el 84 y el 94\,\% del área queda después de él. Ahí el área ordena los entrenamientos por cuánto sobreajustan más que por lo rápido que aprenden, y así se lee.

\paragraph{Mejor \emph{loss} de validación.} El mínimo del \emph{loss} de validación alcanzado
dentro del presupuesto, que mide el mejor nivel intermedio con independencia de la velocidad con que se alcanza. Es una variable de nivel, y desde el 3 de septiembre de 2026 se lee como apoyo del rendimiento final y no de la velocidad, según fija la Sección~\ref{sec:variables}.

\paragraph{Rendimiento final.} El \emph{accuracy} sobre el conjunto de test, evaluado una única
vez al terminar el entrenamiento. Es la variable objetivo certificada, separada de la curva de validación que se monitoriza durante el entrenamiento.

\paragraph{Gap de generalización.} La diferencia final entre el rendimiento en test y en
entrenamiento, medida sobre un subconjunto fijo del conjunto de entrenamiento del mismo tamaño que el de test. En términos de \emph{loss} y de \emph{accuracy},

\begin{equation}
\mathrm{gap}_{\text{loss}} = \mathcal{L}_{\text{test}} - \mathcal{L}_{\text{train}},
\qquad
\mathrm{gap}_{\text{acc}} = \mathrm{acc}_{\text{train}} - \mathrm{acc}_{\text{test}},
\label{eq:gap}
\end{equation}

donde un valor positivo indica sobreajuste. Es la variable objetivo de generalización frente a la que se contrastan, en particular, las métricas que sus trabajos de origen vinculan a este \emph{gap}.

\section{Conclusiones del capítulo}
\label{sec:fund-conclusiones}

Este capítulo ha fijado los dos lados de la correlación que el estudio persigue. Del lado predictor, el gradiente por ejemplo $g_i$ es el objeto común, y las ocho métricas se reducen a menos señales de las que su número sugiere. De las 23 columnas que emiten, 17 son funciones de la matriz de Gram de $256 \times 256$ del \emph{batch} de medición y de sus etiquetas: la m-coherencia, la escala de ruido, la varianza normalizada y la disparidad, porque el gradiente medio de una submuestra es una suma de bloques de esa matriz, y toda la \emph{stiffness} y toda la \emph{gradient confusion}. Fuera quedan GSNR, que agrega momentos por coordenada, y GWA, que compara cada gradiente con los pesos. En población esas 17 columnas se reducen a la razón $\alpha$ entre la señal y la energía del gradiente, a la traza de la covarianza $\operatorname{tr}(\Sigma)$ y a la distribución de los cosenos por pares con su desglose por clase. Las ocho métricas son, por tanto, cinco señales: $\alpha$, $\operatorname{tr}(\Sigma)$, la distribución de cosenos, GSNR y GWA, y GWA es un margen normalizado. Medir las ocho y podar después sigue siendo la decisión correcta, porque las identidades se comprobaron sobre los datos en lugar de suponerse; lo que no cabe es contar cada columna como evidencia independiente, y la Sección~\ref{sec:protocolo-objetivos} fija qué se queda. Del lado a predecir, la eficiencia se mide con seis indicadores leídos sobre la curva de validación y el test final, con el censurado como tratamiento de los entrenamientos lentos que no alcanzan el umbral. Queda pendiente el cómo y el cuándo de la medición: el \emph{batch} de medición fijo, la ventana temporal y los criterios de falsación de cada hipótesis se especifican en el Capítulo~\ref{ch:metodologia}.
